\chapter{Sustainability Plan \& Ethical Considerations}\label{ch:sustain}

\section{Scalability}

\textbf{Current bottlenecks and mitigations:}

\begin{table}[htbp]
\centering
\caption{Scalability analysis}\label{tab:scalability}
\begin{tabularx}{\textwidth}{@{}l X X@{}}
\toprule
\textbf{Component} & \textbf{Potential Bottleneck} & \textbf{Mitigation Strategy}\\
\midrule
Django backend &
  Single-process Daphne struggles under high WebSocket concurrency &
  Run multiple Daphne workers behind Nginx; deploy behind Kubernetes
  HorizontalPodAutoscaler in future\\
PostgreSQL &
  Write throughput for RequestLog ingestion at high volume &
  Partition \texttt{RequestLog} by timestamp; migrate to PostgreSQL
  read replicas for analytics queries\\
VictoriaMetrics &
  Storage node disk exhaustion as more servers are enrolled &
  Cluster architecture (vmstorage+vminsert+vmselect) already supports
  horizontal storage scaling; add vmstorage nodes\\
Celery &
  Task queue depth grows if classify interval is too short &
  Add more Celery workers; tune
  \texttt{REQUEST\_SHIELD\_CLASSIFY\_BATCH\_SIZE}\\
ML microservice &
  CPU-bound inference under large request volumes &
  Deploy multiple ML service replicas; introduce model caching (joblib
  loaded once at startup, already implemented)\\
Redis &
  Memory exhaustion as Celery results accumulate &
  Set \texttt{CELERY\_RESULT\_EXPIRES} TTL; migrate to Redis Cluster\\
\bottomrule
\end{tabularx}
\end{table}

\textbf{Horizontal scaling path:} The stateless nature of the Django backend
and ML microservice allows horizontal scaling via container orchestration
(Kubernetes Deployment with replicas).  VictoriaMetrics cluster already
supports horizontal storage scaling.  A future Kubernetes manifest is
planned for the production deployment.

\section{Flexibility / Customization}

\begin{itemize}
  \item \textbf{Environment-driven configuration:} All tuneable parameters
        (batch size, intervals, Gemini model name, JWT TTLs, VictoriaMetrics
        retention period) are exposed as environment variables in
        \texttt{docker-compose.yml} and loaded via Django settings, requiring
        no code changes to adjust.
  \item \textbf{Modular Django apps:} Each domain (alerts, incidents,
        anomaly detection, Request Shield) is a separate Django app with its
        own models, URLs, and service layer.  New monitoring domains can be
        added without modifying existing apps.
  \item \textbf{Pluggable ML models:} The ML microservice exposes a stable
        \texttt{/train/} and \texttt{/classify/} API; the underlying model
        can be swapped from scikit-learn to XGBoost, LightGBM, or a neural
        network without changing the Django integration code.
  \item \textbf{Feature schema versioning:}
        \texttt{REQUEST\_SHIELD\_ARTIFACT\_FORMAT\_VERSION} ensures
        backward-incompatible feature schema changes are caught at load time.
  \item \textbf{Multi-tenancy:} VictoriaMetrics tenant isolation
        (\texttt{VictoriaMetricsTenant} model) allows each organisation's
        metrics to be physically separated, enabling future per-org retention
        policies or data export.
  \item \textbf{Firebase / notification flexibility:} The
        \texttt{notifications} app abstracts device token registration;
        switching from Firebase to another push provider requires only
        changes to the notification dispatch function.
\end{itemize}

\section{Maintainability and Future Enhancements}

\textbf{Maintenance plan:}
\begin{itemize}
  \item \textbf{Dependency updates:} Monthly review of \texttt{pip-audit}
        and \texttt{flutter pub outdated} output; apply security patches
        within 48 hours.
  \item \textbf{Model retraining:} When admin-labeled production request logs
        exceed 10\% of original CICIDS2017 training set, trigger a retraining
        run via the admin API.
  \item \textbf{Database migrations:} All schema changes via Django
        \texttt{makemigrations}; \texttt{makemigrations --check} run in CI.
  \item \textbf{Documentation:} \texttt{docs/} folder maintained with
        architecture decisions; API changes documented in CHANGELOG.
\end{itemize}

\textbf{Planned future enhancements:}

\begin{enumerate}
  \item \textbf{Kubernetes deployment manifests} -- Helm chart for one-command
        cloud deployment to GKE, EKS, or AKS.
  \item \textbf{Grafana Alloy federation} -- Support multiple Alloy agents
        per server for high-availability metric collection.
  \item \textbf{Log-based anomaly detection} -- Extend the ML pipeline to
        apply anomaly detection to operational log streams, not just metric
        windows.
  \item \textbf{Automated rate limiting} -- Integrate Django REST Framework
        throttling for login endpoints to mitigate brute-force attacks (V-07).
  \item \textbf{Runbook integration} -- Allow organisation admins to upload
        service-specific runbooks that are injected into the Gemini system
        prompt, improving diagnostic relevance.
  \item \textbf{SSO / SAML integration} -- Enterprise single sign-on support
        for organisations with existing identity providers.
  \item \textbf{Mobile push for critical alerts} -- FCM push notifications
        for CRITICAL/HIGH severity alerts, not just incident assignments.
  \item \textbf{UI test automation} -- Flutter integration tests covering
        the registration, dashboard, and chatbot flows.
\end{enumerate}

\section{Ethical Considerations}

\begin{itemize}
  \item \textbf{Data minimisation:} The Request Shield captures HTTP access
        logs including IP addresses, paths, and user agents.  Organisations
        deploying the platform should inform their users of this data
        collection in their privacy policy, in accordance with GDPR and the
        applicable Bangladeshi data protection framework.

  \item \textbf{AI transparency:} The Gemini chatbot's advisory nature is
        disclosed in the UI; responses are framed as suggestions requiring
        human validation.  The system prompt explicitly prohibits the model
        from claiming authority it does not have.

  \item \textbf{Bias and fairness:} The CICIDS2017 dataset reflects 2017
        traffic patterns from a specific lab environment.  The classifier
        may exhibit bias against traffic patterns that were novel in 2017
        but are now common (e.g., HTTP/2, modern browser fingerprints).
        Administrators are advised to review GRAY-zone classification results
        regularly and supplement training data with organisation-specific
        labeled logs.

  \item \textbf{Human oversight:} The Request Shield does \textit{not}
        automatically block IP addresses.  The ThreatSuggestion workflow
        requires explicit admin acceptance before any action is taken,
        preserving human oversight for all security decisions.

  \item \textbf{Environmental responsibility:} The team considered energy
        usage of continuous background ML inference (Celery tasks every
        30--120 seconds).  Interval settings can be tuned to reduce CPU
        usage on constrained hardware.

  \item \textbf{Responsible disclosure:} Any security vulnerabilities
        discovered in the platform should be reported via the GitHub security
        advisory mechanism before public disclosure, following coordinated
        vulnerability disclosure principles.
\end{itemize}

